% --- Άσκηση 38845 (38845.tex) ---
\Askhsh{38845}{4}
\begin{ylh}{2}
    \item 6.1. Η Έννοια της Συνάρτησης
    \item 6.2. Γραφική Παράσταση Συνάρτησης
\end{ylh}
Μια εταιρεία πουλάει φανέλες προς $25$\euro\ τη μία. Το κόστος παραγωγής κάθε φανέλας είναι $5\euro$. Οι πωλήσεις της εταιρείας με αυτή την τιμή $(25)$\euro\ είναι κατά μέσο όρο $100$ φανέλες την ημέρα. Η έρευνα δείχνει ότι για κάθε μείωση της τιμής κατά $1$\euro, η εταιρεία πουλάει $10$ επιπλέον φανέλες.
\begin{alist}
\item Πόσο κέρδος την ημέρα έχει η εταιρεία, όταν πουλάει την κάθε φανέλα με $25$\euro;
\monades{2}
\item Πόσο κέρδος την ημέρα θα έχει η εταιρεία, αν μειώσει την τιμή κάθε φανέλας κατά $2\euro$, δηλαδή από τα $25\euro$ στα $23$\euro;
\monades{5}
\item Αν η εταιρεία μειώσει την τιμή του προϊόντος κατά $x$ \euro, να γράψετε μια σχέση που να εκφράζει το ημερήσιο κέρδος $P$ της εταιρείας ως συνάρτηση του $x$.
\monades{8}
\item Στο παρακάτω σχήμα φαίνεται η καμπύλη της συνάρτησης του κέρδους $P$ σε σχέση με τη μείωση της τιμής $x$.
\begin{rlist}
\item
  Μέχρι πόσο μπορεί να μειώσει η εταιρεία την τιμή κάθε φανέλας, ώστε να μην έχει ζημιά (δηλαδή τα έσοδα να είναι περισσότερα από το κόστος);
\monades{4}
\item
  Να εκτιμήσετε τη μικρότερη τιμή με την οποία μπορεί να πουλάει η εταιρεία την κάθε φανέλα, ώστε το κέρδος της να μην είναι μικρότερο από αυτό που έχει τώρα.
\monades{6}
\end{rlist}
\end{alist}
{\footnotesize Το παραπάνω θέμα αναπτύχθηκε στο πλαίσιο του έργου: «Ανάπτυξη Δοκιμασιών Αξιολόγησης Δεξιοτήτων Εγγραμματισμού στα μαθήματα της Νεοελληνικής Γλώσσας και Λογοτεχνίας, της Άλγεβρας, της Φυσικής και της Χημείας Α' Λυκείου Γενικού Λυκείου» Ανάδοχος: «Ειδικός Λογαριασμός Κονδυλίων Έρευνας (Ε.Λ.Κ.Ε) Πανεπιστημίου Ιωαννίνων» ΑΔΑΜ: 25SYMV016348911 2025-02-20.}
